\subsection{二項係数}

偶数の公式の分数は,階乗で書き直せる.偶数の積からは$2$を1つずつくくり出せるので
\[
 (2m)!!=2^m m!
\]
である.また$(2m)!$は偶数の積と奇数の積に分かれるので
\[
 (2m-1)!!=\frac{(2m)!}{(2m)!!}=\frac{(2m)!}{2^m m!}
\]
となる.これらを式\eqref{式:ウォリス偶数}に代入すると
\begin{equation}\label{式:ウォリスと二項係数}
 I_{2m}=\frac{\pi}{2}\cdot\frac{(2m)!}{2^{2m}(m!)^2}
 =\frac{\pi}{2}\cdot\frac{1}{4^m}\binom{2m}{m}
\end{equation}
を得る.ここで$\binom{2m}{m}=\frac{(2m)!}{(m!)^2}$は二項係数であり,\,$2m$個から$m$個を選ぶ組合せの数${}_{2m}\mathrm{C}_m$に等しい.

二項係数$\binom{2m}{m}$は$m$とともに急速に大きくなる.では,\,$4^m$で割った値はどうなるだろうか.\,$\sqrt m$を掛けて調べてみよう.
\begin{center}
\begin{tabular}{cccccc}
\toprule
$m$ & $1$ & $10$ & $100$ & $1000$ & $10000$\\
\midrule
$\frac{\sqrt m}{4^m}\binom{2m}{m}$ & $0.5$ & $0.5572$ & $0.5635$ & $0.5641$ & $0.5642$\\
\bottomrule
\end{tabular}
\end{center}
値は$0.564$付近へ近づいているように見える.この値を,ウォリスの積分から求めよう.

偶数の公式と奇数の公式を掛けると,分子と分母の偶数がすべて約分されて
\[
 I_{2m}I_{2m+1}
 =\frac{1\cdot3\cdots(2m-1)}{3\cdot5\cdots(2m+1)}\cdot\frac{\pi}{2}
 =\frac{\pi}{2(2m+1)}
\]
となる.これと式\eqref{式:ウォリス隣接比}の$\frac{I_{2m+1}}{I_{2m}}\to1$を組み合わせれば
\begin{align*}
 mI_{2m}^2
 &=m\,I_{2m}I_{2m+1}\cdot\frac{I_{2m}}{I_{2m+1}}\\
 &=\frac{m\pi}{2(2m+1)}\cdot\frac{I_{2m}}{I_{2m+1}}
 \longrightarrow\frac{\pi}{4}
\end{align*}
となる.\,$I_{2m}>0$なので正の平方根を取ることができ,
\[
 \sqrt m\,I_{2m}\longrightarrow\frac{\sqrt\pi}{2}
\]
である.式\eqref{式:ウォリスと二項係数}より$\frac{\sqrt m}{4^m}\binom{2m}{m}=\frac2\pi\sqrt m\,I_{2m}$なので
\begin{equation}\label{式:中央二項係数の極限}
 \lim_{m\to\infty}\frac{\sqrt m}{4^m}\binom{2m}{m}
 =\frac1{\sqrt\pi}
\end{equation}
を得る.表の$0.564$付近の値は,\,$\frac1{\sqrt\pi}=0.56418\ldots$だったのだ.階乗を含む数列の大きさが,三角関数の積分を経由して分かったことになる.

\paragraphbf{硬貨を投げる確率}
硬貨を$2m$回投げて表がちょうど$m$回出る確率は,\,$\frac1{4^m}\binom{2m}{m}$である.式\eqref{式:中央二項係数の極限}によれば,この確率は$m$が大きいとき$\frac{1}{\sqrt{\pi m}}$に近い.例えば$200$回投げて表がちょうど$100$回出る確率は,\,$\frac1{\sqrt{100\pi}}=0.0564\ldots$に近く,実際の値は$0.0563\ldots$である.表と裏が半分ずつになることが最も起こりやすいとはいえ,ちょうど半分になる確率そのものは$0$へ近づいていく.
